% ============================================================
% DROP-IN BLOCK: UPDATED §4.6 + 2×2 ABLATION TABLE
% Replaces the current §4.6 "Structured VAE recovers nonlinear causal variables"
% ============================================================

\subsection{Structured pi-SAE recovers nonlinear causal variables}
\label{sec:vae_results}

The Structured pi-SAE combines two components: (1) a \emph{structured prior}
$p(z_\text{causal} \mid y)$ conditioned on the task label $y$, which pins the
causal direction; and (2) \emph{sparse autoencoding} via a TopK or JumpReLU
activation on $z_\text{causal}$, which enforces a compressed representation.
We evaluate whether both components are necessary via a $2 \times 2$ ablation:
\{structured prior, plain $\mathcal{N}(0,I)$ prior\} $\times$ \{VAE, SAE\}.

Table~\ref{tab:2x2_ablation} reports IIA at $k=2$ across all grokked and
non-grokked operations.

\begin{table}[t]
\centering
\small
\setlength{\tabcolsep}{4pt}
\begin{tabular}{@{}l|cc|cc@{}}
\toprule
& \multicolumn{2}{c|}{\textbf{Plain prior} $\mathcal{N}(0,I)$}
& \multicolumn{2}{c}{\textbf{Structured prior} $p(z \mid y)$} \\
\textbf{Operation} & VAE & SAE & pi-VAE & \textbf{Str.\ pi-SAE} \\
\midrule
\multicolumn{5}{l}{\textit{Grokked operations}} \\
Addition         & 0.02 & 1.00 & 0.00 & \textbf{1.00} \\
Multiplication   & 0.00 & 0.72 & 0.00 & \textbf{1.00} \\
Quartic sum      & 0.02 & 0.32 & 0.00 & \textbf{1.00} \\
IOI (GPT-2)      & 0.56 & 0.83 & 0.95 & \textbf{0.98} \\
\midrule
\multicolumn{5}{l}{\textit{Non-grokked operations (should return $\approx 0$)}} \\
Mixed product    & 0.03 & 0.01 & 0.01 & \textbf{0.04} \\
Squaring         & 0.00 & 0.00 & 0.00 & \textbf{0.00} \\
\bottomrule
\end{tabular}
\caption{$2 \times 2$ ablation: structured prior and sparsity are each
individually insufficient; only their combination is uniformly reliable.
\textbf{Plain SAE}: works on addition (IIA $= 1.00$) but degrades on
multiplication ($0.72$) and quartic sum ($0.32$) because the plain prior
cannot pin the causal direction when the operation's Fourier structure is
more complex.  \textbf{pi-VAE} (structured prior, no sparsity): achieves
$0.00$ on grokked operations---the structured prior provides the right
direction signal, but without sparsity the encoder distributes this
information across all dimensions and cannot be cleanly swapped.
\textbf{Structured pi-SAE}: uniformly $1.00$ on all grokked operations;
$\leq 0.04$ on all non-grokked operations.  Non-grokked operations return
near-zero IIA \emph{despite} the structured prior because the model has no
structured causal variable to recover---confirmed by reconstruction MSE
$\in [11, 24]$ vs.\ $[0.6, 1.2]$ for grokked operations.}
\label{tab:2x2_ablation}
\end{table}

\paragraph{Neither component alone suffices.}
The plain SAE's failure on multiplication and quartic sum is explained by the
structured prior's role: without label conditioning, the sparse encoder has no
target direction and distributes features across the causal and nuisance
dimensions inconsistently across random seeds.  The pi-VAE's failure is
explained by the sparsity's role: a dense causal latent $z_\text{causal}$
requires all $k$ dimensions to be swapped coherently during interchange
interventions, but without sparsity pressure, the encoder spreads the causal
signal across dimensions in a rotation-dependent way that does not survive
swapping.  Together, the structured prior provides the \emph{direction} and
sparsity provides the \emph{compression} needed for faithful interchange.

\paragraph{IOI (GPT-2): nonlinear causal recovery on a pretrained model.}
We validate the Structured pi-SAE on the Indirect Object Identification (IOI)
task~\citep{wang2022interpretability} using GPT-2 Small.  DAS achieves IIA $=
0.30$ at $k=2$, consistent with prior reports that the IOI causal variable
spans multiple attention heads and does not cleanly project onto a
low-dimensional linear subspace.  The Structured pi-SAE achieves IIA $= 0.98$
at $k=2$.  Critically, NL-DAS achieves IIA $= 1.00$---higher than the
structured pi-SAE---but is exposed as vacuous by the diversity ratio:
$\rho_\text{NL-DAS} \approx 0$ vs.\ $\rho_\text{pi-SAE} \approx 0.91$
(Table~\ref{tab:nldas_vacuous}).

The identifiability guarantee (Proposition~\ref{prop:recovery}) applies here
as well: the label-conditional prior satisfies Assumption~\ref{ass:ivae}
because IOI has distinct per-sentence causal labels (the indirect object
identity), the decoder is trained with the ELBO reconstruction constraint
(forcing approximate injectivity), and the supervision term pins the
label-predictive direction.

\paragraph{Reconstruction MSE as a falsifiability criterion.}
The Structured pi-SAE provides a built-in diagnostic for whether a structured
causal variable \emph{exists} to be recovered.  For non-grokked operations,
the reconstruction MSE is $11$--$24$ (vs.\ $0.6$--$1.2$ for grokked
operations), reflecting that the model's activations are high-entropy and
unstructured.  A practitioner applying the pi-SAE to a new task can
interpret high reconstruction MSE with near-zero IIA as evidence that the
model has \emph{no structured causal variable for this operation}, not that
the method failed.  This diagnostic is not available for NL-DAS, which
achieves IIA $= 0.6$--$0.8$ on non-grokked operations (false positives)
precisely because it has no reconstruction constraint.

